% =============================================================
% บทที่ 1 - บทนำ  (ฉบับแก้ไขให้ตรงกับโค้ดใน repo money-printer-research)
% จุดที่มี % TODO คือจุดที่ต้องตรวจ/ใส่ข้อมูลจริงก่อนส่ง
% =============================================================

% -----------------------------------------------------------
\section{ความเป็นมาและความสำคัญของปัญหา}
% -----------------------------------------------------------

ปัจจุบัน ฟุตบอลเป็นกีฬาที่ได้รับความนิยมอย่างมากทั่วโลก และมีการเก็บสถิติการแข่งขันเอาไว้อย่างเป็นระบบ ไม่ว่าจะเป็นผลการแข่งขัน สถิติการครองบอล การยิงประตู ค่าความคาดหวังในการทำประตู (Expected Goals: xG) ไปจนถึงอัตราต่อรองจากเจ้ามือรับพนัน (Odds) โดยเฉพาะการแข่งขันพรีเมียร์ลีกอังกฤษ (English Premier League) ซึ่งมีเตะกันถึง 380 นัดต่อฤดูกาล ทำให้เกิดชุดข้อมูลขนาดใหญ่ที่เหมาะมากสำหรับการนำกระบวนการวิทยาการข้อมูล (Data Science) และการเรียนรู้ของเครื่อง (Machine Learning) มาใช้ศึกษาและวิเคราะห์ความสัมพันธ์ของปัจจัยต่าง ๆ กับผลการแข่งขัน

อย่างไรก็ตาม การจะทำนายผลบอลให้แม่นยำถือเป็นเรื่องที่ท้าทายมาก เพราะผลการแข่งขันขึ้นอยู่กับปัจจัยที่ผันผวนอยู่ตลอดเวลา งานวิจัยในอดีตส่วนใหญ่มักใช้เฉพาะสถิติสกอร์ย้อนหลัง เช่น แบบจำลองเชิงสถิติตระกูล Poisson และ Weibull \cite{dixon1997modelling} หรือระบบการจัดอันดับคะแนนอย่าง Elo Rating, pi-ratings และ Berrar ratings \cite{hvattum2010using, berrar2019incorporating} ซึ่งจากการศึกษาของ Hubáček และคณะ \cite{hubacek2021forty} ที่ทดลองกับชุดข้อมูลกว่า 2.1 แสนนัด จาก 52 ลีกทั่วโลก พบว่าแบบจำลองที่ใช้แค่สถิติย้อนหลังจะทำความแม่นยำได้จำกัดอยู่แค่ช่วงร้อยละ 47.8--48.6 เท่านั้น ในขณะที่ราคาอัตราต่อรองของเจ้ามือรับพนัน ซึ่งประมวลผลจากข้อมูลเชิงลึกและสภาพตลาดจริง กลับทำความแม่นยำได้สูงถึงประมาณร้อยละ 55 และยังไม่มีโมเดลการเรียนรู้ของเครื่องตัวไหนที่ชนะราคาต่อรองได้อย่างชัดเจน

จากจุดนี้ กลุ่มของผู้จัดทำจึงมองเห็นช่องว่างว่า การใช้แค่สถิติย้อนหลังอย่างเดียวอาจไม่เพียงพอ โครงงานนี้จึงมีแนวคิดที่จะดึงข้อมูลมิติอื่น ๆ ที่ไม่ได้อยู่ในสกอร์การแข่งขันเข้ามาช่วยทำนายด้วย ได้แก่ คะแนน Elo Rating เพื่อวัดความแข็งแกร่งของทีม, ดัชนีความเหนื่อยล้าของผู้เล่นตัวจริง, ความหนาแน่นของตารางการแข่งขัน (รวมถึงการแข่งขันถ้วยและรายการยุโรป) และฟอร์มรายบุคคลของผู้เล่นตัวจริงทั้ง 11 คน เช่น ค่า xG ต่อ 90 นาที เพราะงานวิจัยก่อนหน้านี้ชี้ว่า การเอาข้อมูลบริบทเหล่านี้มาร่วมวิเคราะห์กับโมเดลการเรียนรู้ของเครื่อง จะช่วยเพิ่มความแม่นยำในการทำนายได้ดียิ่งขึ้น \cite{berrar2019incorporating}

ดังนั้น โครงงานนี้จึงมุ่งเน้นการรวบรวมข้อมูลการแข่งขันฟุตบอลย้อนหลังตั้งแต่ฤดูกาล 2014--2015 ถึง 2025--2026 จากแหล่งข้อมูลสาธารณะ (Understat และ Transfermarkt) \cite{understat2024, transfermarkt2024} รวมถึงชุดข้อมูลผลการแข่งขันพรีเมียร์ลีกจาก Kaggle \cite{kaggle2024} เพื่อนำมาประมวลผลผ่านท่อส่งข้อมูลที่มีการตรวจสอบโครงสร้างข้อมูล (Schema Validation) สร้างตัวแปรก่อนการแข่งขัน วิเคราะห์ความสำคัญของตัวแปร (Feature Importance) และพัฒนาแบบจำลองสำหรับ\textbf{ทำนายผลการแข่งขัน}ออกเป็น 3 ผลลัพธ์ คือ ทีมเจ้าบ้านชนะ (Home Win) เสมอ (Draw) และทีมเยือนชนะ (Away Win) พร้อมทั้งทดสอบสมมติฐานทางสถิติเกี่ยวกับความได้เปรียบของทีมเจ้าบ้านและความสัมพันธ์ระหว่างคะแนน Elo กับผลต่างประตู เพื่อประโยชน์ในการวิเคราะห์เชิงสถิติ

% -----------------------------------------------------------
\section{วัตถุประสงค์}
% -----------------------------------------------------------

\begin{enumerate}[label=\arabic*)]
    \item เพื่อศึกษาและฝึกปฏิบัติกระบวนการวิทยาการข้อมูลอย่างครบวงจร ตั้งแต่การรวบรวมข้อมูลหลายแหล่ง การเชื่อมโยงข้อมูล การทำความสะอาด การสำรวจข้อมูล ไปจนถึงการพัฒนาและประเมินแบบจำลอง
    \item เพื่อวิเคราะห์หาปัจจัยที่มีอิทธิพลต่อผลการแข่งขันฟุตบอล จากตัวแปรก่อนการแข่งขัน ได้แก่ คะแนน Elo, ดัชนีความเหนื่อยล้า, ความหนาแน่นของตารางแข่งขัน และฟอร์มของผู้เล่นตัวจริง
    \item เพื่อทดสอบสมมติฐานทางสถิติว่า (1) ทีมเจ้าบ้านทำประตูได้มากกว่าทีมเยือน และ (2) ผลต่างคะแนน Elo มีความสัมพันธ์กับผลต่างประตู
    \item เพื่อพัฒนาและประเมินประสิทธิภาพของแบบจำลองในการทำนายผลการแข่งขัน โดยเปรียบเทียบกับค่าพื้นฐาน (Baseline) และผลการศึกษาจากงานวิจัยที่ผ่านมา \cite{hubacek2021forty}
\end{enumerate}

% -----------------------------------------------------------
\section{ขอบเขตและข้อจำกัดของโครงงาน}
% -----------------------------------------------------------

\subsection{ขอบเขตของโครงงาน}

\begin{enumerate}[label=\arabic*)]
    \item \textbf{ขอบเขตด้านข้อมูล:} ใช้ข้อมูลย้อนหลังตั้งแต่ฤดูกาล 2014--2015 ถึง 2025--2026 จากแหล่งข้อมูลสาธารณะดังนี้
    \begin{itemize}
        \item \textit{Understat}: ใช้ดึงตารางการแข่งขันพร้อมจำนวนประตูและค่า xG และสถิติรายบุคคลของผู้เล่นในแต่ละนัด \cite{understat2024}
        \item \textit{Transfermarkt}: ใช้ดึงปฏิทินการแข่งขันเต็มของแต่ละสโมสรทุกรายการ (รวมถึงถ้วยและรายการยุโรป) และแผนการเล่น (Formation) ที่ทีมใช้ \cite{transfermarkt2024}
        \item \textit{Kaggle} (ชุดข้อมูล \textit{marcohuiii/english-premier-league-epl-match-data-2000-2025}): ใช้สำหรับการสำรวจข้อมูลและการทดสอบสมมติฐานเชิงสถิติ กับผลการแข่งขันพรีเมียร์ลีกอังกฤษฤดูกาล 2000--2001 ถึง 2024--2025 \cite{kaggle2024}
    \end{itemize}
    % TODO: FBref ปรากฏเฉพาะใน notebook EDA รุ่นเก่า ถ้าไม่ได้ใช้ในผลงานสุดท้ายไม่ต้องใส่
    \item \textbf{ขอบเขตด้านลีก:} ระบบเก็บรวบรวมข้อมูลและสร้างตัวแปรครอบคลุม 5 ลีกใหญ่ของยุโรป (พรีเมียร์ลีกอังกฤษ, ลาลีกา, บุนเดสลีกา, เซเรียอา และลีกเอิง) ส่วนการทดสอบสมมติฐานเชิงสถิติและแบบจำลอง Elo ใช้ข้อมูลพรีเมียร์ลีกอังกฤษ
    % TODO: ตรวจว่า walk-forward Elo model (/predictions) รันกับลีกไหนจริง แล้วแก้ประโยคนี้ให้ตรง
    \item \textbf{ขอบเขตด้านระบบ:} พัฒนาท่อส่งข้อมูลแบบโมดูลาร์ด้วยภาษา Python ประกอบด้วยการเก็บข้อมูล (\texttt{collect\_data}) การทำความสะอาดและสร้างตัวแปร (\texttt{clean\_data}) การตรวจสอบโครงสร้างข้อมูลด้วย Pandera (\texttt{check\_clean}) การวิเคราะห์ความสำคัญของตัวแปร (\texttt{feature\_importance}) พร้อมชุดทดสอบอัตโนมัติและ API สำหรับเรียกดูผลการทำนาย
    \item \textbf{ขอบเขตด้านตัวแปรเชิงทำนาย (Features):} ใช้เฉพาะข้อมูลที่ทราบก่อนเริ่มการแข่งขัน ได้แก่ คะแนน Elo และผลต่างคะแนน Elo (Elo Difference), ดัชนีความเหนื่อยล้าของผู้เล่นตัวจริง, ความหนาแน่นของตารางแข่งขัน (วันพัก จำนวนนัดใน 7 และ 14 วัน เกมถ้วยหรือเกมยุโรป) และฟอร์มรายบุคคลของผู้เล่นตัวจริง 11 คน
    \item \textbf{ขอบเขตด้านแบบจำลอง:} (1) แบบจำลอง Double Poisson ที่ขับเคลื่อนด้วยผลต่างคะแนน Elo ร่วมกับการแก้ไขแบบ Dixon-Coles \cite{dixon1997modelling} เพื่อคำนวณความน่าจะเป็นของผลการแข่งขัน H/D/A และ (2) แบบจำลอง Gradient Boosting (\texttt{HistGradientBoostingClassifier} จากไลบรารี Scikit-learn \cite{scikit-learn}) เพื่อวิเคราะห์ความสำคัญของตัวแปรด้วยวิธี Permutation Importance โดยใช้ค่าความน่าจะเป็นตามสัดส่วนผลการแข่งขันในอดีต (Class Prior) เป็นเกณฑ์เปรียบเทียบพื้นฐาน
    \item \textbf{ขอบเขตด้านการแบ่งชุดข้อมูล:} แบ่งข้อมูล\textbf{ตามลำดับเวลา (ไม่สุ่ม)} เพื่อป้องกันปัญหาข้อมูลรั่วไหลจากอนาคต (Data Leakage) โดยการวิเคราะห์ความสำคัญของตัวแปรใช้ฤดูกาลก่อนปี 2024 เป็นชุดฝึก และตั้งแต่ปี 2024 เป็นต้นไปเป็นชุดทดสอบ ส่วนแบบจำลอง Elo ใช้การทำนายแบบ Walk-forward คือทำนายแต่ละฤดูกาลด้วยแบบจำลองที่ฝึกจากฤดูกาลก่อนหน้าเท่านั้น และถือว่าฤดูกาล 2014--2015 เป็นช่วงเริ่มต้น (Warm-up) ของคะแนน Elo
    \item \textbf{ตัวชี้วัดประสิทธิภาพ:} ประเมินด้วย Cross-Entropy Loss (Log Loss), Ranked Probability Score (RPS) และ Accuracy พร้อมวัดระดับความสำคัญของตัวแปรด้วย Permutation Importance ส่วนการทดสอบสมมติฐานใช้ Paired t-test, Wilcoxon Signed-rank Test, สัมประสิทธิ์สหสัมพันธ์ของ Pearson และ Spearman
    % TODO: ถ้าจะใส่ F1 / Confusion Matrix / Calibration ต้องเพิ่มโค้ดคำนวณก่อน ไม่งั้นอย่าอ้าง
\end{enumerate}

\subsection{ข้อจำกัดของโครงงาน}

\begin{enumerate}[label=\arabic*)]
    \item ผลการทดสอบสมมติฐานและแบบจำลอง Elo อ้างอิงข้อมูลพรีเมียร์ลีกอังกฤษเป็นหลัก จึงอาจไม่สามารถนำไปใช้ทำนายผลการแข่งขันในลีกอื่นได้ทันที
    \item แบบจำลองทำนายผลลัพธ์ออกเป็น 3 ผลลัพธ์เท่านั้น (เจ้าบ้านชนะ, เสมอ, ทีมเยือนชนะ) และไม่ได้ใช้ราคาอัตราต่อรองของเจ้ามือเป็นตัวแปรหรือเป็นเกณฑ์เปรียบเทียบ จึงไม่สามารถสรุปได้ว่าแบบจำลองเหนือกว่าตลาดการพนัน
    \item ข้อมูลมาจากหลายแหล่งที่เรียกชื่อทีมและชื่อนักเตะไม่เหมือนกัน จึงต้องใช้พจนานุกรมจับคู่ชื่อ (Name Mapping) ซึ่งอาจมีความคลาดเคลื่อนจากการเชื่อมข้อมูล และสโมสรที่จับคู่กับ Transfermarkt ไม่ได้จะมีค่าตัวแปรตารางแข่งขันเป็นค่าว่าง
    \item ข้อมูลของ Understat ครอบคลุมเฉพาะนัดในลีก ดัชนีความเหนื่อยล้ารายผู้เล่นจึงอาจสูงเกินจริงสำหรับสโมสรที่ลงแข่งถ้วยหรือรายการยุโรปด้วย (ชดเชยบางส่วนด้วยตัวแปรตารางแข่งขันระดับทีมจาก Transfermarkt) และตัวแปรบางตัวที่มองไปข้างหน้า เช่น จำนวนวันถึงนัดถัดไป ใช้วันที่แข่งจริง ซึ่งอาจต่างจากที่ทราบก่อนเริ่มนัดในกรณีเลื่อนการแข่งขัน
    \item โครงงานไม่ได้นำข้อมูลการบาดเจ็บหรือปัจจัยเชิงคุณภาพที่ไม่ปรากฏในสถิติ เช่น ขวัญกำลังใจ สภาพอากาศ สภาพสนาม หรือข่าวสารนอกสนาม มาร่วมวิเคราะห์
    \item คะแนน Elo ของฤดูกาล 2014--2015 เริ่มจากค่าเท่ากันทุกทีม จึงมีความน่าเชื่อถือต่ำกว่าฤดูกาลหลัง ๆ
    \item การดำเนินงานอยู่ภายใต้ระยะเวลาจำกัดประมาณ 6 สัปดาห์ และผลลัพธ์ที่ได้มีวัตถุประสงค์เพื่อการศึกษาเชิงสถิติเท่านั้น
\end{enumerate}

% -----------------------------------------------------------
\section{ประโยชน์ที่คาดว่าจะได้รับ}
% -----------------------------------------------------------

\begin{enumerate}[label=\arabic*)]
    \item ผู้จัดทำได้ฝึกปฏิบัติกระบวนการวิทยาการข้อมูลอย่างครบวงจรกับชุดข้อมูลกีฬาจริงที่มีความซับซ้อนและมาจากหลายแหล่ง
    \item ได้ท่อส่งข้อมูลที่ตรวจสอบได้และทำซ้ำได้ (Reproducible) พร้อมแบบจำลองที่ทำนายความน่าจะเป็นของผลการแข่งขันได้ และทราบปัจจัยสำคัญที่มีผลต่อการแข่งขัน
    \item ได้คำตอบเชิงประจักษ์ว่า ตัวแปรกลุ่มใด (Elo, ความเหนื่อยล้า, ตารางแข่งขัน, ฟอร์มผู้เล่น) ให้สัญญาณต่อการทำนายผลการแข่งขันมากที่สุด และแบบจำลองให้ผลดีกว่าค่าพื้นฐานมากน้อยเพียงใด \cite{hubacek2021forty, berrar2019incorporating}
    \item ได้แนวทางและองค์ความรู้ในการนำเทคนิค Data Science และการวิเคราะห์ Feature ไปประยุกต์ใช้กับข้อมูลกีฬาประเภทอื่น ๆ ต่อไป
\end{enumerate}